\begin{afterword}
\summaryhead

Asked what AIs will be like, the author calls it a trick question. Humans form a natural class: as a sexually reproducing species we share one brain architecture. ``AI'', by contrast, names the whole space of possible minds, and two AI designs might differ more than a human and a petunia.

A diagram shows this space as a sphere. Humans are a small dot inside a transhuman region, which lies inside a posthuman region; the rest of the sphere holds minds stranger still, and the whole sphere sits inside a larger space of optimization processes that includes natural selection. Since there are about two to the trillionth power minds that can be specified in a trillion bits, any claim that all minds do something has that many chances to fail, and any claim that some mind does it has that many chances to hold. The usual source of such claims is imagining oneself in the mind's place. The post concludes that one should explain what a mind does from its particular makeup, and not define away minds that behave differently.

\responsehead

The post's logical point is correct. The set of possible minds is far larger than the set of human minds, so a claim about all possible minds needs more support than a claim about humans, and imagining yourself in a mind's place tells you mainly about minds like yours. The rule the post draws, to reason from a particular mind's makeup to what it does, follows. The picture of a space of possible minds is older than the post: Aaron Sloman described the task of mapping it in 1984.

The post uses this to answer the opening question, and its answer is a caution: asking what ``AIs'' will do assumes that they form a natural class, and any two designs ``might'' be less alike than a human and a petunia. Whether the AIs actually built would cluster, as humans do by sharing an architecture, for example because they are made by shared methods, the post does not take up. Quintin Pope later raised this against using the size of mind space as an intuition pump.

In short: a correct warning against generalizing over all possible minds, applied to the question of what AIs will do, without taking up whether the AIs actually built would share a design.
\end{afterword}
